\documentclass[12pt,a4paper]{article}
\usepackage[utf8]{inputenc}
\usepackage[T1]{fontenc}
\usepackage{amsmath,amssymb,amsfonts}
\usepackage{graphicx}
\usepackage{hyperref}
\hypersetup{colorlinks=true,linkcolor=blue,urlcolor=blue}

\title{GRA-Обнуленка: Архитектура гениального AI-агента}
\author{Автор}
\date{\today}

\begin{document}

\maketitle

\begin{abstract}
Представляется архитектура интеллектуального агента, основанная на принципе GRA-Обнуленки — фильтрации состояний по иерархической устойчивости $S$ и энтропии $E$. Агент выбирает действия, максимизирующие произведение $I = S \cdot E$, что обеспечивает баланс между структурной глубиной и творческой новизной. Предлагаются математические основы, алгоритм работы и результаты моделирования, демонстрирующие превосходство над стандартными LLM-подходами в задачах, требующих нетривиального мышления.
\end{abstract}

\section{Введение}

Современные большие языковые модели (LLM) достигают высокой точности за счёт экстремального масштабирования, но часто страдают от отсутствия «интереса» и внутренней мотивации. Мы предлагаем альтернативу: агент, который не просто предсказывает следующий токен, а оценивает каждый возможный ответ по двум измерениям — устойчивости (структурная связность) и энтропии (богатство возможных продолжений). Эта идея восходит к концепции GRA-Обнуленки, где реальность описывается как множество состояний, отфильтрованных по порогу устойчивости.

\section{Математическая модель}

\subsection{Состояние агента}

Пусть $\psi_t$ — состояние агента в момент $t$ (включает контекст, память и внутренние представления). Для каждого состояния определим:

\begin{equation}
S(\psi) = \sum_{k=1}^{N} \lambda_k \log\!\left(1 + \frac{C_k}{V_k}\right),
\end{equation}

где $C_k$ — число устойчивых связей на иерархическом уровне $k$, $V_k$ — общее число узлов, $\lambda_k$ — вес уровня.

Энтропия $E(\psi)$ — мера разнообразия допустимых продолжений (например, энтропия распределения вероятностей по следующим действиям).

\subsection{Фильтр GRA}

Оператор обнуления:

\begin{equation}
\hat{G}[\psi] = \begin{cases}
\psi, & S(\psi) \ge S_{\text{крит}},\\
0, & S(\psi) < S_{\text{крит}}.
\end{cases}
\end{equation}

Состояния с недостаточной устойчивостью исключаются из рассмотрения.

\subsection{Интерес и выбор}

Интерес определяется как:

\begin{equation}
I(\psi) = S(\psi) \cdot E(\psi).
\end{equation}

Агент генерирует набор кандидатов $\{\psi^{(k)}\}$ (ответов, действий) и выбирает:

\begin{equation}
\psi_{t+1} = \arg\max_{\psi^{(k)}: S(\psi^{(k)}) \ge S_{\text{крит}}} I(\psi^{(k)}).
\end{equation}

\subsection{Динамика и сингулярность}

Технологическая сингулярность наступает, когда:

\begin{equation}
\frac{dS}{dt} > \frac{dE}{dt}, \quad \frac{dI}{dt} > 0.
\end{equation}

Агент входит в самоподдерживающийся режим роста сложности и интереса.

\section{Архитектура реализации}

Бэкенд построен на FastAPI и использует модуль `gra_core` для вычисления $S$ и $E$. Генерация кандидатов может выполняться с помощью небольшой языковой модели (например, GPT-2) или эвристических шаблонов. Фронтенд — простой веб-чат, отправляющий запросы к API.

\section{Эксперименты}

Проведена серия экспериментов на следующих задачах:

\begin{itemize}
\item Символьная регрессия (восстановление законов физики по малым данным);
\item Диалоговое взаимодействие (оценка интересности ответов человеком);
\item Креативная генерация текста (сравнение с базовыми LLM).
\end{itemize}

Во всех случаях GRA-агент показал более высокие оценки по «интересности» и структурной связности при сравнимых вычислительных затратах.

\section{Заключение}

Предложенная архитектура гениального AI-агента на основе GRA-Обнуленки обеспечивает естественный баланс между порядком и хаосом. Это открывает путь к созданию систем с внутренней мотивацией, способных к саморазвитию без внешнего учителя. Дальнейшие исследования будут направлены на интеграцию с современными LLM и обучение с подкреплением на основе интереса.

\bibliographystyle{plain}
\begin{thebibliography}{9}
\bibitem{gra-core} GRA-Core: Unified Hierarchical Stability Library, \url{https://github.com/qqewq/GRA-Core-new-Unified-Hierarchical-Stability-Library}
\bibitem{gra-obnulenka} GRA-Обнуленка: грань описываемой реальности, \url{https://github.com/qqewq/GRA-Obnulenka-The-Edge-of-Describable-Reality}
\bibitem{kaplan2020scaling} Kaplan, J. et al. (2020). Scaling Laws for Neural Language Models. arXiv:2001.08361.
\end{thebibliography}

\end{document}